\subsection{Physical-measure parameter estimation}
\label{subsec:p_parameter_estimation}

\subsubsection{SVCJ dynamics and MCMC estimation}
\label{subsubsec:p_svcj_mcmc}

We estimate the physical-measure SVCJ model from daily BTC/USD log returns.
The daily Euler transition is
\begin{align}
r_t &= \mu^P+\sqrt{V_{t-1}}\epsilon^y_t+J_t Z^y_t,
\label{eq:p_svcj_return}\\
V_t &= \alpha^P+\beta^P V_{t-1}
      +\sigma_v^P\sqrt{V_{t-1}}\epsilon^v_t+J_t Z^v_t,
\label{eq:p_svcj_variance}
\end{align}
where $\operatorname{corr}(\epsilon^y_t,\epsilon^v_t)=\rho^P$,
$J_t\sim\operatorname{Bernoulli}(\lambda^P)$,
$Z^v_t\sim\operatorname{Exp}(\mu_v^P)$, and
\begin{equation}
Z^y_t\mid Z^v_t
\sim\mathcal{N}\!\left(\mu_y^P+\rho_j^P Z^v_t,
                        (\sigma_y^P)^2\right).
\end{equation}
The exponential distribution is parameterized by its mean. The restrictions
$\alpha^P>0$ and $0<\beta^P<1$ imply the no-jump long-run variance
$\theta^P=\alpha^P/(1-\beta^P)$ and mean-reversion rate
$\kappa^P=1-\beta^P$.

We use Bayesian data augmentation to update the structural parameters jointly
with the latent variance path, jump indicators, and jump sizes. The first half
of every Markov chain is discarded as burn-in, and posterior means summarize
the retained draws. Successive rolling windows are warm-started from the
preceding posterior mean, while the latent jump indicators remain stochastic
states updated within the sampler rather than externally imposed observations.

\subsubsection{Simulation design for P-parameter recovery}
\label{subsubsec:p_parameter_recovery}

Before estimating the model from observed returns, we evaluate whether the
MCMC implementation can recover known physical-measure parameters and whether
recovery improves with sample length.  The data-generating vector follows the
reference implementation,
\begin{align}
\boldsymbol{\theta}^{P}_0
&=(\mu,\alpha,\beta,\sigma_v,\rho,\lambda,
   \mu_y,\sigma_y,\rho_j,\mu_v,V_0)' \\
&=(0.002,3.5\times10^{-5},0.95,0.005,-0.2,0.05,
   -0.01,0.11,-0.5,0.0006,0.0007)'.
\label{eq:p_simulation_true_parameters}
\end{align}
These values imply $\kappa=0.05$ and $\theta=0.0007$.  Estimation is carried
out in $(\alpha,\beta)$ because this is the sampler's natural parameterization;
the economically interpretable $(\kappa,\theta)$ are calculated from every
retained draw and reported alongside the directly sampled parameters.

We generate 20 independent master paths of 365 daily returns.  Within each
replication, estimation uses the nested prefixes
\begin{equation}
\mathcal{Y}_{30}\subset\mathcal{Y}_{90}
\subset\mathcal{Y}_{180}\subset\mathcal{Y}_{360},
\end{equation}
so changes across $n$ are not confounded by different simulated histories.
For every replication--sample-size pair, four chains start from deliberately
dispersed parameter vectors.  Each chain contains 100,000 MCMC sweeps, of
which the first 50,000 are discarded.  The sampler uses ancillary--sufficiency
interweaving between centered and noncentered variance representations and an
adaptive Fisher-transformed proposal for $\rho$ during burn-in.  Posterior
means and 90\% equal-tail credible intervals are computed from the pooled
retained draws of the four chains.

For parameter $j$ and replication $b$, absolute percentage error is
\begin{equation}
\operatorname{APE}_{jb}(n)
=100\left|\frac{\widehat\theta_{jb}(n)-\theta_{0,j}}
{\theta_{0,j}}\right|.
\label{eq:p_parameter_ape}
\end{equation}
The reported mean APE averages this quantity across the 20 replications.
We additionally calculate across-replication bias, RMSE, and empirical coverage
of the 90\% credible interval.  Mixing is evaluated separately for each
replication--sample-size pair using rank-normalized split $\widehat R$, bulk
effective sample size (ESS), and tail ESS.  A group passes only if every
reported parameter and the stored conditional log-likelihood statistic satisfy
$\widehat R\leq1.01$, bulk ESS $\geq400$, tail ESS $\geq400$, and contain no
nonfinite draws.  Estimates from failed groups remain in the table rather than
being selectively discarded.

\begin{table}[htbp]
\centering
\scriptsize
\resizebox{\textwidth}{!}{%
\begin{tabular}{lrrrrrrrrr}
\toprule
& & \multicolumn{2}{c}{$n=30$} & \multicolumn{2}{c}{$n=90$}
& \multicolumn{2}{c}{$n=180$} & \multicolumn{2}{c}{$n=360$}\\
\cmidrule(lr){3-4}\cmidrule(lr){5-6}\cmidrule(lr){7-8}\cmidrule(lr){9-10}
Parameter & True value & Estimate & APE (\%) & Estimate & APE (\%)
& Estimate & APE (\%) & Estimate & APE (\%)\\
\midrule
$V_0$      & 0.000700 & 0.001716 & 145.2 & 0.001745 & 149.2 & 0.001733 & 147.6 & 0.001594 & 127.8\\
$\mu$      & 0.002000 & 0.001409 & 184.8 & 0.002376 & 140.0 & 0.001358 & 87.3 & 0.001865 & 64.4\\
$\mu_y$    & $-0.010000$ & $-0.030602$ & 443.0 & $-0.006171$ & 480.1 & 0.001230 & 577.9 & 0.010539 & 514.4\\
$\sigma_y$ & 0.110000 & 0.115416 & 7.7 & 0.102549 & 12.9 & 0.095541 & 21.0 & 0.100355 & 21.7\\
$\lambda$  & 0.050000 & 0.050223 & 20.3 & 0.049696 & 32.4 & 0.052934 & 33.8 & 0.047834 & 33.2\\
$\alpha$   & 0.000035 & 0.001018 & 2809.5 & 0.000697 & 1891.9 & 0.000557 & 1491.9 & 0.000231 & 560.5\\
$\beta$    & 0.950000 & 0.140858 & 85.2 & 0.294739 & 69.0 & 0.492150 & 48.2 & 0.774528 & 18.5\\
$\rho$     & $-0.200000$ & 0.022395 & 140.6 & 0.014006 & 133.5 & $-0.135420$ & 116.7 & $-0.175243$ & 88.0\\
$\sigma_v$ & 0.005000 & 0.010086 & 101.7 & 0.008700 & 74.0 & 0.008536 & 70.7 & 0.007985 & 59.7\\
$\rho_j$   & $-0.500000$ & $-0.129865$ & 98.5 & $-0.230553$ & 116.2 & $-0.253463$ & 174.7 & 0.120796 & 208.3\\
$\mu_v$    & 0.000600 & 0.008831 & 1371.8 & 0.005503 & 817.2 & 0.002990 & 398.4 & 0.002004 & 233.9\\
$\kappa$   & 0.050000 & 0.859142 & 1618.3 & 0.705261 & 1310.5 & 0.507850 & 915.7 & 0.225472 & 350.9\\
$\theta$   & 0.000700 & 0.001167 & 75.3 & 0.000960 & 46.4 & 0.001000 & 45.6 & 0.001026 & 47.8\\
\midrule
Convergence pass & & \multicolumn{2}{c}{13/20} & \multicolumn{2}{c}{0/20}
& \multicolumn{2}{c}{0/20} & \multicolumn{2}{c}{1/20}\\
\bottomrule
\end{tabular}%
}
\caption{\textbf{Simulation-based recovery of physical-measure SVCJ parameters.}
The table reports posterior means and mean absolute percentage errors across
20 independent replications.  Each replication uses nested return samples and
four chains with 50,000 retained draws per chain.  The final row reports the
number of replication--sample-size groups satisfying all $\widehat R$ and ESS
criteria.  APE is computed within each replication before averaging and is
therefore not generally equal to the percentage error of the reported mean.}
\label{tab:p_parameter_recovery_simulation}
\end{table}

Table~\ref{tab:p_parameter_recovery_simulation} shows that longer return
histories improve some parameters but do not establish successful joint
recovery.  Mean APE for $\kappa$ declines monotonically from 1,618.3\% at
$n=30$ to 350.9\% at $n=360$, while the directly estimated $\beta$ improves
from 85.2\% to 18.5\%.  Nevertheless, the mean estimate
$\widehat\kappa=0.2255$ at $n=360$ remains more than four times its true value.
The APE of $\theta$ falls from 75.3\% to 46.4\% between $n=30$ and $n=90$ but
does not improve thereafter.  The variance-jump mean $\mu_v$, return-jump mean
$\mu_y$, and co-jump loading $\rho_j$ also remain weakly recovered.

Convergence evidence is more restrictive than the APE pattern.  Thirteen of
20 groups pass all diagnostics at $n=30$, whereas none passes at $n=90$ or
$n=180$ and only one passes at $n=360$.  The higher pass rate for the shortest
sample should not be interpreted as superior identification: its posterior is
broad, so chains can overlap even while estimates remain inaccurate.  These
results indicate that ASIS and adaptive proposals alone do not resolve the
strong posterior dependence and weak finite-sample identification of the full
SVCJ parameter vector.  Accordingly, the exercise is evidence on the current
estimator's limitations rather than validation of rolling estimates.

\subsubsection{Estimation-window selection and QQ diagnostics}
\label{subsubsec:p_window_selection}

The estimation window must balance local adaptation against identification.
A short return window is insufficient for the joint SVCJ specification: it
contains eleven structural parameters and latent states, while the jump block
is informed by relatively few events.  The recovery experiment above shows
that even 360 observations do not reliably identify the full parameter vector.
This concern is consistent with the much longer daily samples used in the
literature. \citet{ErakerJohannesPolson2003}
estimate stochastic-volatility jump models from long S\&P~500 and Nasdaq
histories, and \citet{ChaimLaurini2018} use 1,840 daily Bitcoin observations.

We therefore adopt 1,000 returns as the baseline rolling window. For a market
that trades seven days per week, this corresponds to approximately 2.7 years
and contains about 50 jump days in expectation at the simulation value
$\lambda^P=0.05$.  This choice is a conservative response to the poor recovery
observed at $n\leq360$, not a claim that the present simulation has validated a
1,000-return window.  An expanding window and a 2,000-return window therefore
serve as robustness specifications, and all empirical estimates must be
accompanied by chain-level convergence diagnostics. For target date $t$, the
baseline window contains only returns dated strictly before $t$, which
prevents look-ahead bias.

Model adequacy is assessed with a QQ plot of the standardized diffusion
innovation after removing the posterior expected jump contribution,
\begin{equation}
\widehat\epsilon_t^y
=\frac{r_t-\widehat\mu^P
        -\widehat{J_tZ_t^y}}
       {\sqrt{\widehat V_{t-1}}}.
\label{eq:p_standardized_innovation}
\end{equation}
Under a correctly specified diffusion component, its central quantiles should
be close to those of a standard normal distribution. Tail deviations are
interpreted as residual misspecification, not as a requirement that raw
Bitcoin returns themselves be Gaussian.

\subsection{Risk-neutral parameter estimation}
\label{subsec:q_parameter_estimation}

\subsubsection{Risk-neutral SVCJ specification}
\label{subsec:q_svcj_specification}

We first validate the numerical recovery of the risk-neutral parameters before
applying the calibration procedure to observed option prices. This validation
is important because SVCJ calibration is a nonlinear inverse problem: several
combinations of stochastic-volatility and jump parameters can generate similar
option-price surfaces, Monte Carlo error perturbs the objective function, and a
local least-squares optimizer can converge to different solutions from
different initial values. A satisfactory fit to market prices alone therefore
does not establish that the implementation can recover the parameters that
generated those prices. The simulation study below supplies a controlled
environment in which the data-generating parameter vector is known.

Let $Y_t=\log S_t$. Under the risk-neutral measure $\mathbb{Q}$, the SVCJ
dynamics used in the simulation are
\begin{align}
dY_t
&=\left(r-\delta-\frac{1}{2}V_t-\lambda\kappa_J\right)dt
  +\sqrt{V_t}\,dW_t^S+J_t^Y\,dN_t,
\label{eq:q_svcj_log_price}\\
dV_t
&=\kappa(\theta-V_t)dt
  +\sigma_v\sqrt{V_t}\,dW_t^V+J_t^V\,dN_t,
\label{eq:q_svcj_variance}
\end{align}
where $dW_t^S dW_t^V=\rho\,dt$, $N_t$ is a Poisson process with
intensity $\lambda$, and the simultaneous variance and log-price jumps satisfy
\begin{align}
J_t^V&\sim\operatorname{Exp}(\mu_v),\\
J_t^Y\mid J_t^V&\sim
\mathcal{N}\!\left(\mu_y+\rho_JJ_t^V,\sigma_y^2\right).
\end{align}
Here $\operatorname{Exp}(\mu_v)$ is parameterized by its mean. The
risk-neutral compensator is
\begin{equation}
\kappa_J
=E^{\mathbb{Q}}\!\left[e^{J_t^Y}-1\right]
=\frac{\exp(\mu_y+\tfrac{1}{2}\sigma_y^2)}
       {1-\rho_J\mu_v}-1,
\label{eq:q_svcj_jump_compensator}
\end{equation}
which requires $1-\rho_J\mu_v>0$. This drift adjustment ensures that the
discounted underlying price is a martingale. The specification follows the
affine jump-diffusion framework of \citet{DuffiePanSingleton2000} and the SVCJ
model of \citet{ErakerJohannesPolson2003}.

All parameters and maturities in this diagnostic experiment are expressed in
daily units. With one Euler step per day, the implementation uses full
truncation, $V_k^+=\max(V_k,0)$, and updates
\begin{align}
Y_{k+1}
&=Y_k+\left(r-\delta-\frac{1}{2}V_k^+
            -\lambda\kappa_J\right)
  +\sqrt{V_k^+}\,Z_{k+1}^S+I_{k+1}J_{k+1}^Y,
\label{eq:q_svcj_euler_log_price}\\
V_{k+1}
&=\max\left\{0,
  V_k^++\kappa(\theta-V_k^+)
  +\sigma_v\sqrt{V_k^+}\,Z_{k+1}^V
  +I_{k+1}J_{k+1}^V\right\},
\label{eq:q_svcj_euler_variance}
\end{align}
where
\begin{equation}
Z_{k+1}^V=\rho Z_{k+1}^S
          +\sqrt{1-\rho^2}\,Z_{k+1}^{\perp},
\qquad
I_{k+1}=\mathbf{1}\{U_{k+1}<\lambda\}.
\end{equation}

\subsubsection{Simulation design for Q-parameter recovery}
\label{subsec:q_parameter_recovery_simulation}

The data-generating vector is
\begin{align}
\boldsymbol{\theta}^{\mathbb Q}_0
&=\left(V_0,\kappa,\theta,\sigma_v,\rho,\lambda, \mu_y,\sigma_y,\rho_J,\mu_v\right)^{T} \\ 
&=\left(0.0009,0.7425,0.0008,0.0489,-0.6245,
       0.0033,0.3166,0.1267,-1.0821,0.0206\right)^{T}.
\label{eq:q_simulation_true_parameters}
\end{align}
We set $S_0=40{,}000$ and $r=\delta=0$. The candidate contract grid
contains 52 distinct integer maturities spanning 1 to 365 days and 2,001
equally spaced log-moneyness values from $-0.50$ to $0.50$ at each maturity,
for a total of 104,052 distinct contracts. To avoid placing mechanically
large intrinsic-value observations in the loss, the grid retains an
out-of-the-money put when the strike is below the forward and an
at-the-money or out-of-the-money call otherwise. The contracts are shuffled
once using a fixed seed, and the samples are nested:
\begin{equation}
\mathcal D_{100}\subset\mathcal D_{1000}\subset\mathcal D_{10000}.
\end{equation}
Thus, $n$ denotes the number of distinct option prices used in calibration,
not the number of simulated paths.

For each of five replications, synthetic observed prices are formed from
20,000 paths simulated at $\boldsymbol{\theta}^{\mathbb Q}_0$. Contracts with
the same maturity share the same terminal-price sample, and the discounted
average vanilla payoff gives
\begin{equation}
P_i^{\mathrm{obs}}
=\frac{1}{M_D}\sum_{m=1}^{M_D}
 e^{-rT_i}\max\!\left\{\omega_i
 \left(S_{T_i}^{(m)}-K_i\right),0\right\},
\qquad M_D=20{,}000,
\label{eq:q_simulation_observed_price}
\end{equation}
where $\omega_i=1$ for a call and $\omega_i=-1$ for a put. The synthetic
validation uses standard USD vanilla payoffs so that the Monte Carlo pricing
and optimization pipeline can be isolated from the additional payoff
transformation required by Deribit inverse options. The inverse-payoff
implementation is validated separately before the market-data calibration.

The experiment uses independent common-random-number sets across the data and
estimation stages. Specifically,
\begin{equation}
P_i^{\mathrm{obs}}
=\widehat P_i(\boldsymbol{\theta}^{\mathbb Q}_0;Z^{D}),
\qquad
P_i^{\mathrm{model}}(\boldsymbol{\theta})
=\widehat P_i(\boldsymbol{\theta};Z^{E}),
\qquad Z^{D}\ne Z^{E},
\label{eq:q_independent_crn_design}
\end{equation}
with $M_E=20{,}000$ estimation paths. The estimation draws $Z^E$ are cached
and held fixed at every objective evaluation and for every start within a
replication. Hence the optimizer sees a deterministic objective, but the
true-parameter loss need not be zero because the data and estimation path sets
are independent. This design is more stringent than a same-CRN diagnostic,
under which the true parameter produces zero pricing error by construction.

\subsubsection{Calibration and recovery criterion}
\label{subsec:q_parameter_recovery_estimation}

For every replication and value of $n$, the ten parameters are estimated by
bounded nonlinear least squares using MATLAB's trust-region-reflective
\texttt{lsqnonlin} algorithm. The relative residual for contract $i$ is
\begin{equation}
r_i^{USD}(\boldsymbol{\theta})
=\frac{P_i^{\mathrm{obs, USD}}
       -P_i^{\mathrm{model, USD}}(\boldsymbol{\theta})}
       {\max\{P_i^{\mathrm{obs, USD}},10^{-4}S_0\}},
\label{eq:q_relative_residual_USD}
\end{equation}
or equivalently
\begin{equation}
r_i^{BTC}(\boldsymbol{\theta})
=\frac{P_i^{\mathrm{obs, BTC}}
       -P_i^{\mathrm{model, BTC}}(\boldsymbol{\theta})}
       {\max\{P_i^{\mathrm{obs, BTC}},10^{-4}\}},
\label{eq:q_relative_residual_BTC}
\end{equation}
and the estimator minimizes
\begin{equation}
\widehat{\boldsymbol{\theta}}_n^{\mathbb Q}
=\arg\min_{\boldsymbol{\theta}\in\Theta}
 \sum_{i\in\mathcal D_n}r_i(\boldsymbol{\theta})^2.
\label{eq:q_nonlinear_least_squares}
\end{equation}
The denominator prevents expensive near-the-money contracts from completely
dominating the tail options that help identify jump parameters. The parameter
bounds are
\begin{align*}
\boldsymbol{\ell}
&=(0.00005,0.05,0.00005,0.005,-0.999,
   0.00005,-0.50,0.005,-8.0,0.0005)^{\prime},\\
\boldsymbol{u}
&=(0.01000,2.00,0.01000,0.200,0.500,
   0.05000,0.80,0.500,2.0,0.1000)^{\prime}.
\end{align*}

These bounds are economically and numerically motivated regularization
restrictions rather than parameter ranges copied mechanically from a single
empirical study.  The support restrictions require
$V_0,\theta,\sigma_v,\lambda,\sigma_y,$ and $\mu_v$ to be positive, and
$\rho$ to lie strictly inside the correlation interval.  Because Bitcoin
trades continuously, the daily-variance bounds
$V_0,\theta\in[0.00005,0.01]$ correspond to annualized diffusive
volatilities of approximately 13.5\% to 191\% under a 365-day convention,
which accommodates both tranquil and stressed cryptocurrency regimes.  The
Brownian correlation range $\rho\in[-0.999,0.5]$ permits the empirically
important negative return--variance leverage channel while retaining moderate
positive dependence as a possibility.  The daily jump-probability range
$\lambda\in[0.00005,0.05]$ corresponds to approximately 0.018 to 18.25
expected jump days per year and prevents the jump component from becoming
either exactly absent or observationally indistinguishable from a high-frequency
diffusion.

The remaining jump-size bounds are deliberately broad.  The interval
$\mu_y\in[-0.5,0.8]$ permits large negative and positive conditional log-price
jumps, while $\sigma_y\in[0.005,0.5]$ permits substantial jump-size
dispersion.  The parameter $\rho_J\in[-8,2]$ is a loading of the variance jump
in the conditional mean of the return jump, not a correlation coefficient;
its negative region permits crashes to coincide with upward variance jumps.
Together with $\mu_v\in[0.0005,0.1]$, these bounds guarantee
$1-\rho_J\mu_v>0$ throughout the search region, so the risk-neutral jump
compensator in equation~\eqref{eq:q_svcj_jump_compensator} remains finite.

Finally, $\kappa\in[0.05,2]$ imposes positive variance mean reversion and is
adapted to the daily Euler implementation rather than taken directly from
continuous-time equity-index estimates.  In the shock-free Euler recursion,
the coefficient on lagged variance is $1-\kappa$; setting the upper bound to
two rules out coefficients with absolute value greater than one and therefore
excludes mechanically explosive mean-reversion dynamics.  The endpoint
$\kappa=2$ is retained as a numerical boundary but is not strictly
contractive.  We do not separately impose the diffusion Feller condition
$2\kappa\theta\geq\sigma_v^2$, because the data-generating calibration used
in the recovery experiment need not satisfy it and the implementation instead
enforces nonnegative simulated variance through full truncation.  Boundary-hit
frequencies are reported as an identification diagnostic; systematic contact
with a bound would motivate a dedicated sensitivity calibration rather than an
economic interpretation of the constrained estimate.

We use four deterministic mild-SV initial vectors. They begin from
economically regular stochastic-volatility regimes while holding the jump
parameters at the common neutral values
$(\lambda,\mu_y,\sigma_y,\rho_J,\mu_v)
=(0.003,0,0.05,0,0.02)$. This structured design proved more reliable than
widely dispersed random starts. Parameter scaling is start-dependent,
\begin{equation}
\mathrm{TypicalX}_j=\max\{|\theta_j^{(0)}|,10^{-3}\},
\end{equation}
and forward finite differences use a relative step of $10^{-2}$. The
true-parameter vector is not included among the starts. For each replication
we retain the solution with the smallest relative-residual sum of squares.

Parameter recovery is evaluated using absolute percentage error,
\begin{equation}
\operatorname{APE}_j
=100\left|
\frac{\widehat\theta_j-\theta_{0,j}}
     {\theta_{0,j}}
\right|.
\label{eq:q_parameter_ape}
\end{equation}
Table~\ref{tab:q_parameter_recovery_simulation} reports estimates averaged
across the five replications. APE is computed within each replication and
then averaged, so it is not generally equal to the percentage error of the
reported mean estimate.

\begin{table}[htbp]
\centering
\scriptsize
\resizebox{\textwidth}{!}{%
\begin{tabular}{lrrrrrrr}
\toprule
& & \multicolumn{3}{c}{Mean estimated parameter}
& \multicolumn{3}{c}{Mean APE (\%)}\\
\cmidrule(lr){3-5}\cmidrule(lr){6-8}
Parameter & True value & $n=100$ & $n=1{,}000$ & $n=10{,}000$
& $n=100$ & $n=1{,}000$ & $n=10{,}000$\\
\midrule
$V_0$      & 0.000900 & 0.001038 & 0.000933 & 0.000918 & 29.98 & 9.06  & 3.68\\
$\kappa$   & 0.742500 & 0.396534 & 0.561539 & 0.547842 & 46.59 & 24.37 & 26.22\\
$\theta$   & 0.000800 & 0.000856 & 0.000924 & 0.000897 & 15.13 & 21.49 & 14.86\\
$\sigma_v$ & 0.048900 & 0.038726 & 0.035032 & 0.038401 & 22.93 & 29.50 & 21.47\\
$\rho$     & $-0.624500$ & $-0.288091$ & $-0.667680$ & $-0.614877$ & 53.87 & 18.42 & 17.05\\
$\lambda$  & 0.003300 & 0.003844 & 0.003427 & 0.002961 & 16.63 & 25.50 & 21.96\\
$\mu_y$    & 0.316600 & 0.203652 & 0.318214 & 0.323335 & 35.68 & 23.14 & 14.92\\
$\sigma_y$ & 0.126700 & 0.044407 & 0.120930 & 0.120403 & 64.95 & 11.17 & 10.10\\
$\rho_J$   & $-1.082100$ & 1.848960 & $-0.145671$ & 1.054524 & 270.87 & 262.33 & 197.45\\
$\mu_v$    & 0.020600 & 0.018796 & 0.012498 & 0.009348 & 23.28 & 42.39 & 54.62\\
\midrule
Mean APE   & & & & & 57.99 & 46.74 & 38.23\\
Median APE & & & & & 32.83 & 23.75 & 19.26\\
\bottomrule
\end{tabular}%
}
\caption{\textbf{Simulation-based recovery of risk-neutral SVCJ parameters.}
The table reports results from five independent replications. Synthetic
observed prices and model prices use separate sets of 20,000 Monte Carlo paths.
For each replication, the calibration selects the lowest-loss solution among
four deterministic mild-SV starts. The option samples are nested across
$n=100$, $1{,}000$, and $10{,}000$. Mean APE is calculated parameter by
parameter across replications. The final two rows summarize the ten reported
parameter-specific APE values.}
\label{tab:q_parameter_recovery_simulation}
\end{table}

The cross-sectional increase in option count improves recovery on average.
The mean parameter APE decreases from 57.99\% at $n=100$ to 46.74\% at
$n=1{,}000$ and 38.23\% at $n=10{,}000$; the corresponding median falls from
32.83\% to 23.75\% and 19.26\%. Recovery improves especially for $V_0$,
$\rho$, $\mu_y$, and $\sigma_y$. The improvement is not uniform, however.
The co-jump parameter $\rho_J$ remains weakly identified, and the APE of
$\mu_v$ increases with $n$. These patterns are consistent with compensation
among jump intensity, jump-size, and co-jump parameters. They also show why
pricing fit alone cannot be treated as evidence of parameter identification.
Nevertheless, the decline in the central tendency of APE demonstrates that a
richer option cross-section provides useful recovery information, while the
remaining jump-parameter instability motivates the multiple-start,
diagnostic, and robustness procedures used in the market-data calibration.
